\documentclass[11pt]{article}
% Кодировка и язык
\usepackage[utf8]{inputenc}           % Кодировка исходного файла (UTF-8)
\usepackage[T2A]{fontenc}             % Кодировка шрифтов для кириллицы
\usepackage[english,russian]{babel}   % Локализация: русские названия, переносы

% Поиск в PDF
\usepackage{cmap}                     

% Поля страницы
\usepackage{geometry}
\geometry{left=2cm, right=2cm, top=2cm, bottom=2cm}

% Математика
\usepackage{amsmath, amssymb, amsthm} % Базовые пакеты для математики
\usepackage{mathtools}                % Расширение amsmath (доп. символы и окружения)

% Картинки
\usepackage{graphicx}                 % Вставка картинок (\includegraphics)
\graphicspath{{pictures/}}            % Путь к папке с картинками
\usepackage{float}                    % Позволяет использовать [H] для жёсткого позиционирования
\usepackage{wrapfig}                  % Обтекание текста вокруг картинок

% Таблицы
\usepackage{array}                    % Расширенные возможности для таблиц
\usepackage{booktabs}                 % Красивые таблицы (\toprule, \midrule, \bottomrule)
\usepackage{multirow}                 % Объединение строк в таблицах
\usepackage{tabularx}                 % Таблицы с автоматической шириной колонок
\usepackage{ragged2e}                 % Улучшенное выравнивание (используется в tabularx)

% Списки
\usepackage{enumitem}                 % Гибкая настройка списков (itemize, enumerate)

% Колонтитулы
\usepackage{fancyhdr}                 

% Гиперссылки
\usepackage{xcolor}                   % Работа с цветами
\usepackage{hyperref}                 % Гиперссылки, оглавление, перекрёстные ссылки

% Умные ссылки
\usepackage{cleveref}                 

% Типографика
\usepackage[expansion=false]{microtype}                % Улучшает переносы и межсловные расстояния

% Подписи к картинкам и таблицам
\usepackage{caption}                  % Настройка подписей (шрифт, отступы)

% Рисование схем и графиков
\usepackage{tikz}                     % Рисование схем, графиков, диаграмм
\usepackage{pgfplots}                 % Построение графиков функций (на основе tikz)
\pgfplotsset{compat=newest}           % Версия совместимости pgfplots

% Алгоритмы и код
\usepackage{algorithm}                % Окружения для алгоритмов
\usepackage{algorithmic}              % Синтаксис для алгоритмов
\usepackage{listings}                 % Вставка кода с подсветкой синтаксиса


% Настройка нумерации секций и подсекций
\renewcommand{\thesection}{\arabic{section}}                   
\renewcommand{\thesubsection}{\thesection.\arabic{subsection}} 


% Настройка списков: убираем лишние отступы слева
\setlist[itemize]{leftmargin=*}
\setlist[enumerate]{leftmargin=*}


% Определение окружений для теорем, лемм, определений и т.д.
% Нумерация привязана к номеру секции
\newtheorem{theorem}{Теорема}[section]
\newtheorem{definition}{Определение}[section]
\newtheorem{remark}{Замечание}[section]
\newtheorem{lemma}{Лемма}[section]
\newtheorem{example}{Пример}[section]
\newtheorem{consequence}{Следствие}[section]
\newtheorem{algorithmm}{Алгоритм}
\newtheorem{question}{Вопрос}[section]
\newtheorem{answer}{Ответ}[section]
\newtheorem{statement}{Утверждение}[section]


% Определение цветов
\definecolor{pblue}{rgb}{0.13,0.13,1}    
\definecolor{pgreen}{rgb}{0,0.5,0}       
\definecolor{pred}{rgb}{0.9,0,0}         
\definecolor{pgrey}{rgb}{0.46,0.45,0.48}

\newcommand{\lectureNumber}{7}
\newcommand{\lectureTitle}{Алгоритм $\varepsilon$-аппроксимации}
\newcommand{\courseName}{Формальные языки и сложность вычислений}
\newcommand{\semester}{3 семестр, 2026/2027 уч. год}

\pagestyle{fancy}
\fancyhf{}
\fancyhead[L]{Не очень-то вежливо уходить из гостей сразу, как только ты наелся}
\fancyhead[R]{\thepage}
\setlength{\headheight}{15pt}

\hypersetup{
    colorlinks=true,
    linkcolor=blue!70!black,
    citecolor=blue!70!black,
    filecolor=blue!70!black,
    urlcolor=blue!70!black,
}

\begin{document}


\begin{titlepage}
    \centering
    
    {\large Московский физико-технический институт \par}
    {\large Высшая школа программной инженерии \par}
    \vspace{2cm}

    {\Large \textbf{\courseName} \par}
    {\large \semester \par}
    
    \vspace*{\fill}

    {\Huge \textbf{Лекция \lectureNumber} \par}
    {\Huge \textbf{\lectureTitle} \par}
    
    \vspace*{\fill}
    
    {\large Лектор: Зухба Анастасия Викторовна \par}
    {\large Автор: \href{https://t.me/mallina_olga}{Винни-пух затехал для всех-всех-всех} \par}

\end{titlepage}

\pagenumbering{arabic}

\section{Алгоритм $\varepsilon$-аппроксимации}

Способы формулировки NP-полных задач:
\begin{enumerate}
    \item Разрешение языка (то есть рассматривается какой-то язык и какое-то слово, вопрос: лежит ли данное слово в языке? $x \overset{?}{\in} L$)
    \item Задача поиска\\
    Если NP-полная задача --- это вопрос о существовании, то соответствующая задача поиска --- указать требуемый объект. 
    \item Задача оптимизации --- это такая задача, что для каждого её экземпляра $x$ мы имеем множество решений $F(x)$ и для каждого решения $s\in F(x)$ определена функция стоимости $c(s)$, являющаяся положительным целым числом. Оптимальная стоимость решения тогда определяется как
    \[ OPT(x) = \underset{s\in F(x)}{max} c(s) \]  
    (или min, в зависимости от задачи)
\end{enumerate}

\begin{statement}
    Если задача NP-полная и есть оракул, разрешающий её за один шаг, то существует полиномиальный алгоритм, решающий задачу поиска и использующий оракулы. Это называется самосводимостью. 
\end{statement}

\begin{proof}
        Строим самосводимость для 3-SAT: предположим, что у нас есть оракул, который по произвольной формуле $\phi$ может определить есть ли выполняющий набор.
        Тогда наш алгоритм в таком случае будет действовать следующим образом: 
        \begin{enumerate}
            \item Формула, которую мы рассматриваем равна исходной: $\psi = \phi$
            \item Проверим выполнимость $\psi$ 
            \item Пока есть параметры у формулы $\psi$, подставим $\psi(x_i=1)$. Если выполнимо, то $\psi = \psi(x_i=1)$, а $x_i=1$ добавим в выполняющий набор.
            Иначе $\psi = \psi(x_i=0)$, а в выполняющий набор добавим $x_i=0$.   
        \end{enumerate}
\end{proof}

\begin{definition}
    (что такое МТ с оракулом) У нас есть какая-то  МТ. У неё есть какая-то дополнительная лента. В любой момент времени мы можем вывести на дополнительную ленту какое-нибудь слово, затем переходим в специальное состояние, обращаясь к оракулу, и после этого МТ переводит либо в состояние $q_{yes}$, либо в состояние $q_{no}$. 
\end{definition}


\begin{definition}
    (Алгоритм $\varepsilon$-аппроксимации) Пусть есть алгоритм $M$, который по экземпляру задачи $x$ возвращает решение $M(x)\in F(x)$. Тогда мы будем говорить,
    что $M$ --- \textbf{это алгоритм $\varepsilon$-аппроксимации} (где $\varepsilon$ между 0 и 1 включительно), если для любого $x$:
    \[ \frac{|c(M(x))-OPT(x)|}{max \{ c(M(x)), OPT(x) \} } \leq \varepsilon \] 
\end{definition}

\begin{definition}
    \textbf{Порог аппроксимации $A$} --- это точная нижняя грань всех $\varepsilon>0$, таких что существует алгоритм $\varepsilon$- аппроксимации для $A$, работающий за полиномиальное время. 
\end{definition}

\begin{example}
    Рассмотрим задачу поиска для вершинного покрытия $\mathrm{MINVERTEXCOVER}$.В качестве $x$ имеем некоторый граф: $x=G$. 
    Тогда $\mathrm{OPT(x)} = $ ''настоящее'' минимальное вершинное покрытие.
    Тогда, если мы предъявим какой-то алгоритм $c(M(x))$, строящий вершинное покрытие (то, что мы построим обязано быть вершинным покрытием, так как это и есть принадлежность множеству $F(x))$, то
    $c(M(x)) \geq OPT(x)$. Докажем, что порог аппроксимации для данной задачи не выше $\frac{1}{2}$. Для этого предъявим алгоритм с $\frac{1}{2}$. \\
    То есть, мы требуем:
    \[ \frac{c(M(x))- OPT(x)}{c(M(x))} \leq \frac{1}{2} \]   
    Это значит: 
    \[ c(M(x)) \leq 2\ OPT(x) \]   
    Пусть покрытие $C=\{ \} .$ Пока в $G$ есть рёбра:
    \begin{enumerate}
        \item возьмём ребро $(u,v)$
        \item добавим $u$, $v$ к покрытию $C$
        \item удалим $u$, $v$ из графа $G$ со всеми смежными рёбрами
    \end{enumerate}
    $C($разбитое на пары естественным образом$)$ образует паросочетание. 
    Но если есть некоторое паросочетание, то любое вершинное покрытие, в том числе оптимальное, обязано содержать хотя бы по одной вершине из паросочетания для каждого ребра. Значит
    \[ OPT(x) \ge \frac{1}{2} |C| \] 

\end{example}


\begin{example}
    Но есть задачи, которые плохо приближаются. Например, $TSP$.  Условие задачи коммивояжера:
        \begin{enumerate}
        \item у нас есть полный взвешенный граф $G$
        \item $W$ --- набор весов рёбер
        
    \end{enumerate}
    Небоходимо найти путь, проходящий через все вершины графа ровно по одному разу (причём старт и финиш совпадают) наименьшей длины.
    Пусть у нас есть полиномиальный по $|x|$ алгоритм $\varepsilon$-аппроксимации $TSP$ с $\varepsilon\ge 1$.
    Идея: заметим, что оптимальный путь (если он есть) нам показывает гамильтонов путь. \\
    От противного: пусть существует алгоритм $\varepsilon$-аппроксимации с $\varepsilon<1$. Этот алгоритм полиномиален по $|x|. $ 
    Покажем, как с его помощью решить задачу разрешения языка $\mathrm{HAMCYCLE}$. 
    Зададим веса в графе $G$:
    \[ \forall e\in E \ \omega_e =1 \] 
    \[ \forall e\notin E \ \omega_e = M, \quad M = \frac{|V|}{1- \varepsilon} \] 
    Если $G\in \mathrm{HAMCYCLE}$, то $OPT(x) = |V|$.\\
    Если $G\notin \mathrm{HAMCYCLE}$, то $OPT(x) \ge |V|-1+ M $ (но так как M мы выбираем сами и можем выбрать его достаточно большим, то сразу можем сказать: $OPT(x)\ge M)$
    Имеем:
    \[ \frac{c(M(x))-OPT(x)}{c(M(x))} = 1 - \frac{OPT(x)}{c(M(x))} \le \varepsilon \]   
    Отсюда получим:
    \[ c(M(x)) \le \frac{OPT(x)}{1-\varepsilon} \]   
    Тогда, если мы получили решение $ c(M(x)) > \frac{|V|}{1-\varepsilon}$, поскольку $1-\varepsilon$ положительно, в итоге получим, что $OPT(x)> |V| $. А значит, цикла нет. 

\end{example}




\subsection{Полиномиальная аппроксимационная схема}


\begin{definition}
    \textbf{Полиномиальная аппроксимационная схема}  для задачи оптимизации А это алгоритм, такой что для любого подаваемого
    $\varepsilon>0$ и экзамепляра задачи х он является алгоритмом $\varepsilon$-аппроксимации, и работает за полином от $|x|$ (полином может зависеть от $\varepsilon$)
\end{definition}

\begin{definition}
    \textbf{Полностью полиномиальная аппроксимационная схема} --- это полиномиальная аппроксимационная схема, которая работает и за полином от $\frac{1}{\varepsilon}$
\end{definition}




























\end{document}